\section{Historia material}

La historiografía popular de la nanotecnología se fundamenta en un mito de
origen: la noción de que la disciplina surgió de la conferencia dictada por
Richard P. Feynman en 1959 \cite{feynmanTheresPlentyRoom2011}.
No obstante, el registro histórico revela que su influencia en la génesis del
campo fue nula. Aunque Feynman articuló experimentos mentales sobre la
miniaturización, su discurso carecía de los formalismos matemáticos y los
protocolos experimentales necesarios para guiar la manipulación atómica.

El análisis bibliométrico de Chris Toumey evidencia la desconexión entre
esta profecía y la práctica científica \cite{toumeyReadingFeynmanNanotechnology2008}.
La conferencia de 1959 generó apenas tres citas en la década de 1960 y cuatro en
la de 1970. El texto fue ignorado por la comunidad científica hasta ser
resucitado décadas después. Autores como K. Eric Drexler utilizaron a Feynman
como argumento de autoridad retrospectivo, buscando conferir el prestigio de un
laureado a agendas tecnológicas que carecían de rigor analítico.

La verdadera génesis de la nanotecnología reside en los avances formales de
la física de la materia condensada y la mecánica cuántica. La
manipulación atómica fue inaugurada por hitos instrumentales empíricos,
destacando la invención del microscopio de efecto túnel en 1981 por Gerd
Binnig y Heinrich Rohrer, y la posterior manipulación de átomos de
xenón por Don Eigler. Binnig y Rohrer no se fundamentaron en el discurso de
Feynman. Su desarrollo respondió estrictamente a la necesidad de resolver
anomalías en la física de superficies bajo vacío ultraalto.

De manera análoga, el descubrimiento de los fullerenos en 1985 por Kroto,
Smalley y Curl respondió a interrogantes fundamentales en la espectroscopía
molecular y la astroquímica. La base material de la disciplina se
construyó sobre el escrutinio de interacciones físicas. El progreso dependió
invariablemente del estudio de la naturaleza, demostrando que la innovación
tecnológica es consecuencia directa de la ciencia básica, no de premisas
teleológicas orientadas a la aplicación inmediata.
